\section*{第 \ref{ch:b2:cell-differentiation} 章——细胞分化：骨骼肌细胞}

\begin{solution}{exo:b2:cell-differentiation:1}
\omterm{def:b2:cell-differentiation:differentiation}{分化}：某一细胞类型基因与功能的稳定表达，
决定：命运在表达之前即已确定，
\omterm{def:b2:cell-differentiation:differentiation}{潜能}：尚未关闭的命运范围。合子：
全能；\omterm{def:b2:mammal-reproduction:implantation}{内细胞团}细胞：多能；卫星细胞：
单能（仅限肌肉）——但仍是\omterm{def:b2:cell-differentiation:differentiation}{干细胞}；\omterm{def:b2:cell-differentiation:myogenesis}{肌纤维}：
已\omterm{def:b2:cell-differentiation:differentiation}{分化}，有丝分裂后，无\omterm{def:b2:cell-differentiation:differentiation}{潜能}。
\end{solution}

\begin{solution}{exo:b2:cell-differentiation:2}
皮\omterm{def:b2:cell-differentiation:fibre}{肌节}细胞（Pax3）$\to$ 诱导的\omterm{def:b2:cell-differentiation:myogenesis}{成肌细胞}（\omterm{def:b2:cell-differentiation:myogenesis}{MyoD}、Myf5），
在 FGF 存在期间持续分裂 $\to$ 退出细胞周期，肌细胞生成素 $\to$ 排列
与融合形成\omterm{def:b2:cell-differentiation:myogenesis}{肌管} $\to$ 肌肉基因开启，\omterm{def:b2:cell-differentiation:fibre}{肌原纤维}
装配（MRF4）$\to$ 成熟\omterm{def:b2:cell-differentiation:myogenesis}{肌纤维}，具周边核，另有
静止的卫星细胞（Pax7）。
\end{solution}

\begin{solution}{exo:b2:cell-differentiation:3}
该实验表明，已\omterm{def:b2:cell-differentiation:differentiation}{分化}细胞的细胞核保留了
构建整个动物所需的全部基因，卵细胞质
能够重置其程序：\omterm{def:b2:cell-differentiation:differentiation}{分化}是可逆的，而非 DNA 丢失。
但它并未证明每个已\omterm{def:b2:cell-differentiation:differentiation}{分化}的细胞核都能被重置（多数
移植失败），亦未说明重置如何实现，也未证明
已\omterm{def:b2:cell-differentiation:differentiation}{分化}细胞本身能够就地改变命运。
\end{solution}

\begin{solution}{exo:b2:cell-differentiation:4}
Z 盘界定\omterm{def:b2:cell-differentiation:fibre}{肌节}；细肌丝锚定于
其上；粗肌丝以 M 线为中心；A 带
（粗肌丝）与 I 带（仅细肌丝）。缩短时，细肌丝
向 M 线滑动，Z 盘彼此靠近：I 带与 H 区
变窄，A 带——即粗肌丝的长度——
不变，两种肌丝都不缩短。
\end{solution}

\begin{solution}{exo:b2:cell-differentiation:5}
$0.2/2.5\times 10^{-6} = \num{80000}$ 个\omterm{def:b2:cell-differentiation:fibre}{肌节}；$0.2/25\times 10^{-6}
= \num{8000}$ 个细胞核；每个细胞核对应一个\omterm{def:b2:cell-differentiation:myogenesis}{成肌细胞}，故约 \num{8000}
个\omterm{def:b2:cell-differentiation:myogenesis}{成肌细胞}发生了融合（卫星细胞后来还会增加）。
\end{solution}

\begin{solution}{exo:b2:cell-differentiation:6}
FGF 诱导产生 Id，Id 结合 \omterm{def:b2:cell-differentiation:myogenesis}{MyoD} 所需的
伙伴蛋白（形成 DNA 结合二聚体所需）；\omterm{def:b2:cell-differentiation:myogenesis}{MyoD} 虽存在却闲置，细胞持续
分裂。撤去 FGF 后 Id 降解，\omterm{def:b2:cell-differentiation:myogenesis}{MyoD} 形成二聚体，肌细胞生成素
被激活，细胞退出细胞周期并发生融合。\omterm{def:b2:cell-differentiation:myogenesis}{成肌细胞}
若持续产生 Id，则永不融合：它们无限分裂，
不形成肌肉。
\end{solution}

\begin{solution}{exo:b2:cell-differentiation:7}
条件：$0.4m = m^{2}/(1 + m^{2}) + 0.02$。当 $m = 2.1$ 时：左边
$0.84$，右边 $4.41/5.41 + 0.02 = 0.835$：稳态。低
态：对于小的 $m$，$0.4m \approx m^{2} + 0.02$，故 $m \approx
0.055$（左边 $0.022$，右边 $0.023$）。升至 1，高于
约 $0.41$ 的阈值，细胞升至 $2.1$ 并停留：它已
处于已决定的状态。
\end{solution}

\begin{solution}{exo:b2:cell-differentiation:8}
取 $k = 0.6$：当 $m = 1$ 时，降解量 $0.6$ 超过合成量
$0.52$；当 $m = 0.5$ 时，$0.30 > 0.22$；当 $m = 2$ 时，$1.2 > 0.82$：降解线
在低态之外处处位于曲线之上，而低态
是唯一的稳态。若 \omterm{def:b2:cell-differentiation:myogenesis}{MyoD} 降解过快，
细胞便无法维持开启态：它会回复到未\omterm{def:b2:cell-differentiation:differentiation}{分化}状态，
无论如何诱导——均不形成肌肉。
\end{solution}

\begin{solution}{exo:b2:cell-differentiation:9}
(a) 无\omterm{def:b2:cell-differentiation:myogenesis}{成肌细胞}，亦无骨骼肌：两种因子对决定作用
互为冗余，因此两者都须缺失。(b) \omterm{def:b2:cell-differentiation:myogenesis}{成肌细胞}形成并排列，
但从不融合，也不合成肌肉蛋白：肌细胞生成素是
\omterm{def:b2:cell-differentiation:differentiation}{分化}所必需的。(c) 基本正常，因为 Myf5 可替代
\omterm{def:b2:cell-differentiation:myogenesis}{MyoD} 参与决定。
\end{solution}

\begin{solution}{exo:b2:cell-differentiation:10}
在成纤维细胞中，肌肉基因位于 \omterm{def:b2:cell-differentiation:myogenesis}{MyoD} 能够打开的染色质中；
在肝细胞中，它们可能被包裹于异染色质中，或
肝脏自身的\omterm{prop:b2:cell-differentiation:transcriptional}{主调控因子}可能抑制它们，因此该细胞不具有
\omterm{def:b2:vertebrate-development:induction}{感受态}。检验：先用使染色质松散的药物处理肝细胞，
（组蛋白去乙酰化酶抑制剂）后再加入 \omterm{def:b2:cell-differentiation:myogenesis}{MyoD}，或
将 \omterm{def:b2:cell-differentiation:myogenesis}{MyoD} 与可替代肝脏程序的因子一同加入，
并统计肌球蛋白的表达。
\end{solution}

\begin{solution}{exo:b2:cell-differentiation:11}
持续的低频放电使胞内钙离子长时间升高，
钙激活的通路开启了慢肌球蛋白、线粒体生成、
肌红蛋白及毛细血管生长的基因，
并关闭快肌亚型。\omterm{def:b2:cell-differentiation:myogenesis}{肌纤维}的主调控程序
（\omterm{def:b2:cell-differentiation:myogenesis}{MyoD} 状态、\omterm{def:b2:cell-differentiation:fibre}{肌节}结构）未受影响：
训练改变的是哪些肌肉基因被表达，而非
细胞的身份，神经发出的任何信号都无法重新开启
在决定时已关闭的软骨程序。
\end{solution}

\begin{solution}{exo:b2:cell-differentiation:12}
自激活因子的开启态是反馈环路的一个稳定点，
由持续合成维持，而非依赖基因本身的任何改变；
染色质标记另加一层机制，使沉默的
基因保持关闭。因此，重置需将回路
压到阈值以下，并同时重新打开染色质，
卵细胞质或四种因子只能在一小部分
细胞中完成：这是可能的，因为没有任何东西丢失，但效率低下，因为
两把相互独立的锁须同时开启。
\end{solution}

\begin{solution}{pb:b2:cell-differentiation:1}
\textbf{1.} $0.2/2.5\times 10^{-6} = \num{80000}$ 个\omterm{def:b2:cell-differentiation:fibre}{肌节}；
$0.2/25\times 10^{-6} = \num{8000}$ 个细胞核。
\textbf{2.} 每根\omterm{def:b2:cell-differentiation:myogenesis}{肌纤维}约 \num{8000} 个；整块
肌肉为 $4\times 10^{9}$ 个。
\textbf{3.} $\pi(30\times 10^{-6})^{2}\times 0.2 = \qty{5.7e-10}{m^3}$
每根\omterm{def:b2:cell-differentiation:myogenesis}{肌纤维}；整块
肌肉为 $2.8\times 10^{-4}$ \unit{m^3}，\qty{0.28}{L}。
\textbf{4.} $5.7\times 10^{-10}/8000 = \qty{7e-14}{m^3} =
\qty{70000}{\micro m^3}$：每个细胞核相当于 35 个普通细胞的体积。
\textbf{5.} \omterm{def:b2:cell-differentiation:myogenesis}{肌纤维}内部充满了\omterm{def:b2:cell-differentiation:fibre}{肌原纤维}，它们必须
从头到尾连续不断；位于边缘的细胞核使晶格
保持连续，并靠近膜及其信号。
\textbf{6.} $4\times 10^{9}/2000 = 2\times 10^{6}$：需 21 次倍增，耗时 10.5
天。
\textbf{7.} $0.4m = m^{2}/(1 + m^{2}) + 0.02$。当 $m = 0.055$ 时，左侧
为 $0.022$，右侧 $0.003 + 0.02 = 0.023$。
\textbf{8.} 当 $m = 2.1$ 时，左侧 $0.84$，右侧 $0.815 + 0.02 = 0.835$。
\textbf{9.} 当 $m = 0.41$ 时，左侧 $0.164$，右侧 $0.144 + 0.02 = 0.164$。
在其略低处，合成小于降解，故 $m$ 下降
趋向 $0.055$；略高处，合成大于降解，$m$
上升趋向 $2.1$：任何偏移都会增大——不稳定。
\textbf{10.} $0.6 > 0.41$：第一个细胞进入开启态并
\omterm{def:b2:cell-differentiation:differentiation}{分化}；$0.3 < 0.41$：第二个细胞回落至 $0.055$ 并
不\omterm{def:b2:cell-differentiation:differentiation}{分化}。
\textbf{11.} 取 $b = 0.4$ 时，合成曲线从 $0.4$ 出发，
始终位于直线 $0.4m$ 之上，直至 $m \approx 3.3$：低处的交点
消失了。细胞无需诱导即可自发\omterm{def:b2:cell-differentiation:differentiation}{分化}。
\textbf{12.} 分裂使 $m$ 减半；子细胞继承 $2.1/2 =
1.05$，高于阈值 $0.41$，并回升至 $2.1$：该状态
在每个子细胞中重新建立。只要减半后的水平仍高于阈值，
定向便得以保留。
\textbf{13.} 每毫米 $4$ 个 $\times$ \qty{200}{mm} $= 800$ 个卫星细胞，
每根\omterm{def:b2:cell-differentiation:myogenesis}{肌纤维}；\num{8000} 个细胞核中的 \qty{10}{\%}，即 \num{800} 个，须
予以替换。
\textbf{14.} 受损段含 $80$ 个卫星细胞；$800/80 =
10$，故需 4 次倍增（$\times 16$，得 \num{1280}：800 个融合，
480 个用于恢复细胞库）——\qty{72}{h}，即三天。
\textbf{15.} 分裂仅需三天，而观察到的修复需两至三周。其余时间用于
巨噬细胞清除碎片、激活延迟、融合、
\omterm{def:b2:cell-differentiation:fibre}{肌原纤维}组装、神经再支配及
新段成熟。
\textbf{16.} \omterm{def:b2:cell-differentiation:myogenesis}{MyoD}（及肌细胞生成素）须将肌肉基因重新开启
于被激活的细胞中；分裂与\omterm{def:b2:cell-differentiation:differentiation}{分化}互斥
——Id 必须下降、细胞周期停止，然后肌细胞生成素才能发挥作用，
并发生融合。
\textbf{17.} $0.95^{n} = 0.5$：$n = 13.5$，约十四次损伤。
\textbf{18.} $800\times 0.95^{20} = 800\times 0.36 \approx 290$.
\textbf{19.} 每次收缩都会撕裂\omterm{def:b2:cell-differentiation:myogenesis}{肌纤维}，每次撕裂都会动用
细胞库，而细胞库每次都缩小几个百分点；经多年后，它
耗尽，被撕裂的\omterm{def:b2:cell-differentiation:myogenesis}{肌纤维}不再重建，成纤维细胞
以胶原蛋白填补空隙。
\textbf{20.} $(80/60)^{2} = 1.78$；细胞核 $8000\times 1.78 = \num{14200}$：
约多出 \num{6200} 个，来自卫星细胞。
\textbf{21.} $0.28\times 1.78 = \qty{0.5}{L}$.
\textbf{22.} 发生改变的有：肌球蛋白同工型基因、线粒体与
毛细血管密度、肌红蛋白。未改变的有：肌肉身份（\omterm{def:b2:cell-differentiation:myogenesis}{MyoD}
状态）、\omterm{def:b2:cell-differentiation:fibre}{肌节}结构、细胞核及其基因组。
\textbf{23.} \omterm{def:b2:cell-differentiation:myogenesis}{肌纤维}为有丝分裂后细胞，不能增殖；唯一途径是
增粗\omterm{def:b2:cell-differentiation:myogenesis}{肌纤维}，以及借卫星细胞使每根\omterm{def:b2:cell-differentiation:myogenesis}{肌纤维}
获得更多细胞核。
\textbf{24.} 训练：仅能有限增粗，因为\omterm{def:b2:cell-differentiation:myogenesis}{肌纤维}无法增加
细胞核；损伤：无法修复——受损段由
纤维状组织替代，如同肌营养不良中的情形。
\textbf{25.} 每根\omterm{def:b2:cell-differentiation:myogenesis}{肌纤维} \num{8000} 个细胞核；阈值 $m \approx 0.41$；
修复一次约需卫星细胞分裂三天。
\end{solution}
